\documentclass[11pt]{article}
\usepackage[a4paper,margin=27mm]{geometry}
\usepackage{amsmath,amssymb,amsthm}
\usepackage[T1]{fontenc}
\usepackage{hyperref}
\newtheorem{proposition}{Proposition}
\title{An atomic Puiseux monoid of integer rank one}
\author{ziangni-sys}
\date{Conjecture 00000002208}
\begin{document}
\maketitle
\noindent\textbf{Claim addressed.} Both versions of the source assert that a monoid generated by $1$ and positive real numbers is atomic if and only if the integer rank of its generators is at least $2$. The necessity direction is false.

\begin{proposition}
The additive monoid $M=\langle 1,2\rangle\subseteq\mathbb Q_{\geq0}$ is atomic, but the $\mathbb Z$-module generated by its generators has rank $1$.
\end{proposition}
\begin{proof}
By definition,
\[
 M=\{a+2b:a,b\in\mathbb N\}=\mathbb N.
\]
Here $\mathbb N$ includes zero. Both specified generators are strictly positive; the generator $1$ is present as required. This is also a Puiseux monoid under the usual convention of an additive submonoid of nonnegative rational numbers. Embedding $\mathbb Q$ in $\mathbb R$ gives exactly the same example under the source's positive-real convention.

The only additive unit of $M$ is zero: if $a,b\in M$ and $a+b=0$, nonnegativity implies $a=b=0$. An atom in this reduced additive monoid is therefore a nonzero element $u$ for which every decomposition $u=a+b$ with $a,b\in M$ has $a=0$ or $b=0$. The element $1$ is an atom, since two natural numbers summing to $1$ cannot both be positive. Every $n\in M$ is a finite sum of $n$ copies of this atom. For $n=0$ this is the empty sum. Thus $M$ is atomic. In fact, any $n\geq2$ decomposes as $1+(n-1)$ into two nonzero elements, so $1$ is its only atom.

The integer span of the actual generators is
\[
 L=\operatorname{span}_{\mathbb Z}\{1,2\}
   =\mathbb Z\cdot 1\subseteq\mathbb Q,
\]
because $2=2\cdot1$ and $1$ is among the generators. The map $\mathbb Z\to L$, $k\mapsto k\cdot1$, is a bijective $\mathbb Z$-linear map: it is surjective by the description of $L$ and injective because $1\ne0$ in $\mathbb Q$. Consequently $L$ is a free module of rank $1$, which is strictly less than $2$. This proves the counterexample.
\end{proof}

\noindent\textbf{Formal correspondence.} The accompanying Lean project defines $M$ as the actual additive-submonoid closure of $\{1,2\}$ in $\mathbb Q$, proves that its carrier consists exactly of natural-number casts, and proves the zero-unit property. Its atom definition uses decompositions by elements of this closure. The atomicity theorem supplies a finite indexed list of atoms whose sum is each given monoid element, including zero. The rank theorem concerns Mathlib's actual \texttt{Module.rank} of the actual integer span. It is obtained from an explicit linear equivalence with $\mathbb Z$, rather than an assigned numerical rank.

\medskip
\noindent No restriction in either source excludes redundant positive generators or requires density. Even the singleton generating family $\{1\}$ would yield the same rank-one atomic monoid. This submission disproves the stated universal criterion; it makes no claim about a criterion with additional hypotheses.
\end{document}
